SELinux uses various terms to distinguish kinds of components during runtime.
\begin{itemize}
  \item The \emph{object} of an action, e.g.\ a file being read, has an assigned \emph{type}.
  \item The object also has a \emph{class}, e.g.\ file, directory, socket, or process.
  \item The \emph{subject} of an action, e.g.\ the program, runs with a \emph{domain} type.
  \item When loaded, objects and subjects have a security \emph{context}, which includes a \emph{user}, \emph{role}, \emph{type}, and optionally a \emph{sensitivity} and a \emph{category}.
    \begin{itemize}
        \item Users are not standard Linux users, but are types assigned to users that cannot be overridden, e.g.\ \verb|staff_u| and \verb|system_u|.
        \item Users have roles, e.g.\ \verb|staff_u| may have the \verb|sysadm_r| administrator role.
      \item Sensitivities and categories are used for clearance levels in Multi-Level/Category Security schemes, which are not relevant for our work.
    \end{itemize}
  \item \emph{Labels} assigned to files on the filesystem indicate a context, e.g. a file's type or, indirectly, a program's domain.
  \item Domains may \emph{transition} to other domains based on policy rules.
  \item Types may be associated with type \emph{attributes}, which act as an inheritance mechanism.
\end{itemize}

There may be multiple SELinux policy implementations that can accomplish the same rules.
This nuance is not discussed in most basic SELinux policy tutorials,
largely because SELinux policies for most applications do not need to be as strict as ours and follow a relatively standard template.

SELinux has three distinct formats for its policy: a high-level language (HLL), the Common Intermediate Language (CIL), and binary.
The binary format is primarily used for compiled runtime policies.
The HLL and CIL formats are intended to be human-readable and compiled to the binary format.
The standard SELinux tooling converts the HLL to CIL for compilation;
this makes it easy to combine HLL policy with CIL policy manually.

The main factor influencing the complexity of SELinux policy is the Reference Policy \cite{SELinuxProject:2026:ReferencePolicy}.
The Reference Policy is written in the HLL and serves as the basis for most SELinux policies.
It includes some widely useful macros that set up commonly necessary rules,
and it includes conventional configuration for things like administrative users and common Linux programs.
In the Reference Policy, \emph{unconfined} contexts are essentially unrestricted by SELinux.
We address this lack of restrictions in our custom policy (\zcref[S]{sec:implementation-selinux}).

A specific example of a macro included in the Reference Policy is \verb|files_type(x)|, which nominally associates the type \verb|x| with the \verb|file_type| attribute.
The Reference Policy then further includes permissions allowing unconfined contexts to access and manipulate files with \verb|file_type| label \verb|x|.
This is reasonable in most situations where an administrator should be easily allowed to manipulate files.
However, this is not good for our work, as unconfined contexts are given the following extensive set of permissions:
\begin{lstlisting}[style=selinux-policy]
allow files_unconfined_type file_type:file { exec_file_perms manage_file_perms mounton quotaon relabel_file_perms watch watch_mount watch_reads watch_sb watch_with_perm };
\end{lstlisting}
% TODO: cite https://github.com/SELinuxProject/refpolicy/blob/2cba8023863718709d0349faf62a9f4da2248a3f/policy/modules/kernel/files.te#L230,
% TODO: cite https://github.com/SELinuxProject/refpolicy/blob/2cba8023863718709d0349faf62a9f4da2248a3f/policy/support/obj_perm_sets.spt#L158
We found ourselves concerned with
\verb|exec_file_perms|, which is a macro that includes \verb|open| and \verb|read|,
and with \verb|relabel_file_perms|, which includes \verb|relabelfrom| and \verb|relabelto|.

Another macro in the Reference Policy is \verb|unconfined_run_to(domain_t, execfile_t)|,
which allows \verb|unconfined_t| domains to transition to \verb|domain_t| by executing a file with label \verb|execfile_t|.
It produces the following rules:
\begin{lstlisting}[style=selinux-policy]
type_transition unconfined_t execfile_t : process domain_t;
allow unconfined_t domain_t : process { transition };
role unconfined_r types domain_t;
userdom_use_user_terminals(domain_t)
\end{lstlisting}

In Fedora's \verb|targeted| policy, which is based on the Reference Policy, % TODO: cite
unconfined contexts are also able to transition to any domain type and inspect the memory of any program. % TODO: improve

After parsing through these highly-permissive macro definitions in the Reference Policy, it is possible to recreate them while excluding undesirable permissions,
or modify them to replicate their behavior for other types.
However, this can be inconvenient since these details are hidden in the source of the Reference Policy.

% TODO : cite addition of deny rules
The recent addition of \emph{deny rules} in CIL makes a positive step in this regard.
Deny rules explicitly restrict permissions and are capable of overriding existing policy.
They can easily compensate for pre-existing SELinux policy modules which leave unintended security holes,
like policies that use the previously mentioned macros.
This makes deny rules a suitable alternative to manually recreating existing policies with minor modifications,
as shown in \zcref[S]{sec:implementation-selinux}.
% TODO : flesh this out more clearly
% The \verb|typeattribute|/\verb|typeattributeset| declaration is necessary to determine which types should be denied in the last deny rule
% via logical expressions like \verb|not|, \verb|and|, \verb|or|, etc.
% This is because, in CIL policy, rules (\verb|allow|, \verb|deny|, etc.) must take an identifier in the first and second arguments.
% In our policy, this allows us to deny reads from every domain except \domaint.

Some sources refer to \emph{neverallow} rules and incorrectly imply they behave similarly to \emph{deny} rules.
In reality, neverallow rules are compile-time checks that are supposed to trigger warnings upon violation by allow rules.
However, sometimes they get missed depending on the compilation setup and no warning is produced.